\subsection{Exact dynamic reduction and finite-time perturbation selection}
\label{subsec:dynamic_finite_time_operator}

This subsection closes the theoretical gap between the frozen-time quadratic
descriptor in Eq.~\eqref{eq:dynamic_descriptor} and an observable perturbation
amplified over a finite, evolving loading history.  The construction is for the
single-crystal HCP base state.  Experimental parameter identification and
independent full-field material validation are outside the present scope.

\subsubsection{Admissible base history and operator hypotheses}

Let $I=[t_0,t_f]$, let $k>0$, and let $\boldsymbol n\in\mathbb S^2$.  The
following hypotheses define the domain on which the reduction is asserted.
\begin{enumerate}
\item[(H1)] The homogeneous base history is continuous and piecewise $C^1$,
  remains on one recorded smooth constitutive branch on $I$, and supplies
  continuous tangent blocks (piecewise continuity suffices for propagation,
  but not for every threshold-continuity statement below).
\item[(H2)] The density and heat capacity are strictly positive, the continuous
  tangent blocks are finite and bounded, and the mechanical inertia and
  differential storage blocks retain constant rank on $I$.
\item[(H3)] The 18-by-18 algebraic microfield block is uniformly invertible:
  $\inf_{t\in I}\sigma_{\min}
  (\boldsymbol K_{\zeta\zeta}(t;k,\boldsymbol n))>0$ in the locked assembled
  physical coordinates for every admitted $(k,\boldsymbol n)$; the infimum is
  also uniform over each compact parameter domain used below.
\item[(H4)] The eight-coordinate $\mathrm{SL}(3)$ chart is imposed before a
  spectrum or propagator is interpreted.  The continuous local symbol does not
  impose finite-cell wavevector quantization; rigid translation is projected
  only in the $k=0$ limit.  A periodic finite cell instead supplies a separate
  discrete Bloch/Floquet wavevector set and its own gauge treatment.
\item[(H5)] State scales, input and output selectors, the gain threshold, the
  direction--wavenumber search domain, and all numerical tolerances are fixed
  before the resulting optimum is inspected.
\end{enumerate}
If H1 fails at a switching surface, an ordinary Jacobian does not define the
claimed operator.  If H3 fails, the Schur reduction below is disabled and the
augmented descriptor is retained as a possible algebraic or microforce
bifurcation.  Thus invertibility is an admission condition, not a conclusion
manufactured by regularization.

\subsubsection{The 84-coordinate quadratic pencil}

For the locked ordering
\begin{equation}
 \widehat{\boldsymbol z}_{84}
 =
 \left[
  \widehat{\boldsymbol u}_3,
  \widehat T,
  \widehat{\boldsymbol\zeta}_{18},
  \widehat{\boldsymbol a}_{62}
 \right]^{\mathsf T},
 \label{eq:dyn_84_state}
\end{equation}
write the dynamic Fourier operator as the quadratic matrix polynomial
\begin{equation}
 \boldsymbol{\mathcal Q}
 (s;t,k,\boldsymbol n)
 =
 \boldsymbol K(t;k,\boldsymbol n)
 +s\boldsymbol E(t)
 +s^2\boldsymbol M(t),
 \qquad
 \boldsymbol{\mathcal Q}\widehat{\boldsymbol z}_{84}=\boldsymbol0.
 \label{eq:dyn_quadratic_pencil}
\end{equation}
The only nonzero block of $\boldsymbol M$ is the 3-by-3 mechanical block
$\boldsymbol R$; the nonzero differential block of $\boldsymbol E$ is
\begin{equation}
 \boldsymbol C
 =\operatorname{diag}(\rho_0c,\boldsymbol I_{62}).
 \label{eq:dyn_mass_storage_blocks}
\end{equation}
No artificial inertia or first-order storage is assigned to
$\boldsymbol\zeta$.  All $\mathrm{i}k$ and $k^2$ terms, including the
directional acoustic, conduction, and gradient-recovery terms, remain in
$\boldsymbol K(t;k,\boldsymbol n)$.

Partition
\begin{equation}
 \widehat{\boldsymbol y}
 =
 [\widehat{\boldsymbol u}_3,\widehat{\boldsymbol d}_{63}]^{\mathsf T},
 \qquad
 \widehat{\boldsymbol d}
 =
 [\widehat T,\widehat{\boldsymbol a}_{62}]^{\mathsf T},
 \label{eq:dyn_retained_state}
\end{equation}
and place $\widehat{\boldsymbol\zeta}$ second in the block partition.  Because
its rows and columns carry no time derivative,
\begin{equation}
 \boldsymbol{\mathcal Q}
 =
 \begin{bmatrix}
  \boldsymbol K_{yy}+s\boldsymbol E_{yy}+s^2\boldsymbol M_{yy}
    & \boldsymbol K_{y\zeta}\\
  \boldsymbol K_{\zeta y} & \boldsymbol K_{\zeta\zeta}
 \end{bmatrix}.
 \label{eq:dyn_partitioned_pencil}
\end{equation}

\paragraph{Proposition 1 (exact algebraic reduction).}
Under H3, the algebraic amplitude and the retained stiffness are
\begin{align}
 \widehat{\boldsymbol\zeta}
 &=-\boldsymbol K_{\zeta\zeta}^{-1}
       \boldsymbol K_{\zeta y}\widehat{\boldsymbol y},
 \label{eq:dyn_zeta_reconstruction}\\
 \boldsymbol K_{\zeta\mathrm{-red}}^{(66)}
 &=\boldsymbol K_{yy}
   -\boldsymbol K_{y\zeta}\boldsymbol K_{\zeta\zeta}^{-1}
    \boldsymbol K_{\zeta y}.
 \label{eq:dyn_effective_stiffness}
\end{align}
Moreover,
\begin{equation}
 \det\boldsymbol{\mathcal Q}
 =\det\boldsymbol K_{\zeta\zeta}\,
  \det\!\left(
    \boldsymbol K_{\zeta\mathrm{-red}}^{(66)}
    +s\boldsymbol E_{yy}+s^2\boldsymbol M_{yy}
  \right),
 \label{eq:dyn_determinant_factorization}
\end{equation}
so the augmented and reduced problems have the same finite eigenvalues,
including algebraic multiplicity, and every retained eigenvector reconstructs
a unique 84-coordinate amplitude through Eq.~\eqref{eq:dyn_zeta_reconstruction}.

\emph{Proof.}
The lower block row of Eq.~\eqref{eq:dyn_partitioned_pencil} gives
Eq.~\eqref{eq:dyn_zeta_reconstruction}.  Substitution gives
Eq.~\eqref{eq:dyn_effective_stiffness}; block Gaussian elimination gives
Eq.~\eqref{eq:dyn_determinant_factorization}.  The eliminated block is
independent of $s$, so the operation neither creates nor deletes a finite root.
The complete block calculation and the converse reconstruction are recorded in
Appendix~\ref{app:operator_blocks}. \hfill$\square$

\subsubsection{Nonsingular 69-state generator}

Write $\boldsymbol K_{\zeta\mathrm{-red}}^{(66)}$ in the
$[\widehat{\boldsymbol u},\widehat{\boldsymbol d}]$ ordering as
$\boldsymbol K_{uu}$, $\boldsymbol K_{ud}$,
$\boldsymbol K_{du}$, and $\boldsymbol K_{dd}$.  The reduced time-domain
equations are
\begin{align}
 \boldsymbol R\,\ddot{\widehat{\boldsymbol u}}
 +\boldsymbol K_{uu}\widehat{\boldsymbol u}
 +\boldsymbol K_{ud}\widehat{\boldsymbol d}&=\boldsymbol0,
 \label{eq:dyn_reduced_momentum}\\
 \boldsymbol C\,\dot{\widehat{\boldsymbol d}}
 +\boldsymbol K_{du}\widehat{\boldsymbol u}
 +\boldsymbol K_{dd}\widehat{\boldsymbol d}&=\boldsymbol0.
 \label{eq:dyn_reduced_storage}
\end{align}
Introduce physical velocity
$\widehat{\boldsymbol v}=\dot{\widehat{\boldsymbol u}}$ and the state
\begin{equation}
 \widehat{\boldsymbol q}
 =
 [\widehat{\boldsymbol u}_3,
  \widehat{\boldsymbol v}_3,
  \widehat{\boldsymbol d}_{63}]^{\mathsf T}\in\mathbb C^{69}.
 \label{eq:dyn_generator_state}
\end{equation}
Then
\begin{equation}
 \dot{\widehat{\boldsymbol q}}
 =\boldsymbol A_{\mathrm{dyn}}(t;k,\boldsymbol n)
  \widehat{\boldsymbol q},
 \qquad
 \boldsymbol A_{\mathrm{dyn}}
 =
 \begin{bmatrix}
  \boldsymbol0&\boldsymbol I&\boldsymbol0\\
  -\boldsymbol R^{-1}\boldsymbol K_{uu}&\boldsymbol0&
  -\boldsymbol R^{-1}\boldsymbol K_{ud}\\
  -\boldsymbol C^{-1}\boldsymbol K_{du}&\boldsymbol0&
  -\boldsymbol C^{-1}\boldsymbol K_{dd}
 \end{bmatrix}.
 \label{eq:dyn_generator_matrix}
\end{equation}

\paragraph{Proposition 2 (finite-spectrum equivalence).}
Under H2--H3, $s$ is a finite root of
Eq.~\eqref{eq:dyn_quadratic_pencil} if and only if it is an eigenvalue of
$\boldsymbol A_{\mathrm{dyn}}$.  Hence the admitted dynamic pencil has 69
finite roots, counted with algebraic multiplicity: six mechanical first-order
coordinates and 63 thermal/internal-state coordinates.

More precisely, with
$\boldsymbol{\mathcal Q}_{\rm eff}=
\boldsymbol K_{\zeta\mathrm{-red}}^{(66)}
+s\boldsymbol E_{yy}+s^2\boldsymbol M_{yy}$,
\begin{equation}
 \det(s\boldsymbol I_{69}-\boldsymbol A_{\rm dyn})
 =\det(\boldsymbol R^{-1})\det(\boldsymbol C^{-1})
  \det\boldsymbol{\mathcal Q}_{\rm eff}(s).
 \label{eq:dyn_generator_determinant}
\end{equation}

\emph{Proof.}
For a modal solution $\widehat{\boldsymbol q}(t)=
\boldsymbol q_s\exp(st)$, the first block row of
Eq.~\eqref{eq:dyn_generator_matrix} gives
$\widehat{\boldsymbol v}=s\widehat{\boldsymbol u}$.  The remaining rows reduce
exactly to Eqs.~\eqref{eq:dyn_reduced_momentum} and
\eqref{eq:dyn_reduced_storage}; Proposition~1 supplies the unique algebraic
amplitude.  Conversely, a finite pencil eigenvector and
$\widehat{\boldsymbol v}=s\widehat{\boldsymbol u}$ satisfy every generator
row.  Nonsingularity of $\boldsymbol R$ and $\boldsymbol C$ excludes a hidden
differential constraint. \hfill$\square$

This result concerns finite roots.  The 18 eliminated microfield coordinates
remain structural algebraic coordinates in a first-order augmented descriptor;
they are not reinterpreted as very fast physical modes.

\subsubsection{Nonautonomous finite-time propagator}

For fixed $(k,\boldsymbol n)$, H1--H3 make
$\boldsymbol A_{\mathrm{dyn}}(\cdot;k,\boldsymbol n)$ piecewise continuous.
The initial-value problem therefore has the unique absolutely continuous
solution
\begin{equation}
 \widehat{\boldsymbol q}(t)
 =\boldsymbol\Phi(t,t_0;k,\boldsymbol n)
  \widehat{\boldsymbol q}(t_0),
 \qquad
 \dot{\boldsymbol\Phi}
 =\boldsymbol A_{\mathrm{dyn}}\boldsymbol\Phi,
 \qquad
 \boldsymbol\Phi(t_0,t_0)=\boldsymbol I_{69}.
 \label{eq:dyn_propagator_ivp}
\end{equation}
Equivalently,
\begin{equation}
 \boldsymbol\Phi(t,t_0;k,\boldsymbol n)
 =\mathcal T\exp\!\left[
   \int_{t_0}^{t}\boldsymbol A_{\mathrm{dyn}}
   (\xi;k,\boldsymbol n)\,\mathrm d\xi
  \right],
 \label{eq:dyn_time_ordered_exponential}
\end{equation}
where $\mathcal T$ denotes chronological ordering.  Equation
\eqref{eq:dyn_time_ordered_exponential}, rather than an integral of the frozen
spectral abscissa, is the primary finite-time object.  It retains mode rotation,
branch interaction, and non-normal transient amplification.

At every time, the eliminated field is reconstructed from the simultaneous
retained state using Eq.~\eqref{eq:dyn_zeta_reconstruction}.  No
$\dot{\boldsymbol K}_{\zeta\zeta}$ connection term appears in the 69-state
equations because $\boldsymbol\zeta$ has no differential row or column; such a
term would arise only if a time derivative of the reconstructed algebraic
observable itself were requested.

\subsubsection{Fixed physical scaling and observable gain}

Because displacement, velocity, temperature, plastic-chart coordinates, and
dislocation densities have different units, an unscaled Euclidean singular
value is not a physical invariant.  Fix one positive diagonal scale
$\boldsymbol D_q$ for the entire history and define
\begin{equation}
 \widehat{\boldsymbol q}=\boldsymbol D_q\widetilde{\boldsymbol q},
 \qquad
 \widetilde{\boldsymbol A}_{\mathrm{dyn}}
 =\boldsymbol D_q^{-1}\boldsymbol A_{\mathrm{dyn}}\boldsymbol D_q,
 \qquad
 \widetilde{\boldsymbol\Phi}
 =\boldsymbol D_q^{-1}\boldsymbol\Phi\boldsymbol D_q.
 \label{eq:dyn_scaled_generator}
\end{equation}
A time-dependent $\boldsymbol D_q$ is not used; it would add
$-\boldsymbol D_q^{-1}\dot{\boldsymbol D}_q$ to the scaled generator and
would change the observable being measured.

Let $\boldsymbol P_{\rm in}$ and $\boldsymbol P_{\rm out}$ be preregistered
real coordinate-selection matrices.  The pointwise observable gain is
\begin{equation}
 \mathcal G_{D}(t;k,\boldsymbol n)
 =\sigma_1\!\left[
   \boldsymbol P_{\rm out}\widetilde{\boldsymbol\Phi}
   (t,t_0;k,\boldsymbol n)\boldsymbol P_{\rm in}^{\mathsf T}
  \right].
 \label{eq:dyn_observable_gain}
\end{equation}
The localization observable used here selects the 63 temperature and active
crystal-state coordinates at input and output and excludes displacement and
velocity from the singular-value objective.  Elastic inertia nevertheless
remains fully coupled inside $\boldsymbol A_{\mathrm{dyn}}$.

Two gains must be distinguished. The projected quantity
$\mathcal G_D$ in Eq.~\eqref{eq:dyn_observable_gain} answers the conditioned
localization input--output question and is used for direction--wavenumber
selection. The like-for-like comparison with the spectral abscissa instead
uses the complete reduced state,
\begin{equation}
 \mathcal G_D^{\rm full}(t;k,\boldsymbol n)
 =\|\widetilde{\boldsymbol\Phi}(t,t_0;k,\boldsymbol n)\|_2,
 \qquad
 \boldsymbol P_{\rm in}=\boldsymbol P_{\rm out}=\boldsymbol I_{69}.
 \label{eq:dyn_full_state_gain}
\end{equation}
This separation is necessary because the selected 63-coordinate constitutive
subspace is not invariant under the HCP generator: mechanical states can feed
the selected output and selected inputs can generate mechanical response.
Consequently, a projected gain cannot in general be assigned the full-state
frozen spectral abscissa without introducing a controllability/observability
mismatch.

\paragraph{Proposition 3 (equivalent norms and conditioned onset).}
Let $\boldsymbol D_1$ and $\boldsymbol D_2$ be positive diagonal scales and
let the same coordinate selectors be used.  Define
\begin{align}
 c_+^{12}
 &=\left\|\boldsymbol P_{\rm out}\boldsymbol D_1^{-1}
   \boldsymbol D_2\boldsymbol P_{\rm out}^{\mathsf T}\right\|_2
   \left\|\boldsymbol P_{\rm in}\boldsymbol D_2^{-1}
   \boldsymbol D_1\boldsymbol P_{\rm in}^{\mathsf T}\right\|_2,\\
 c_-^{12}
 &=\left[
   \left\|\boldsymbol P_{\rm out}\boldsymbol D_2^{-1}
   \boldsymbol D_1\boldsymbol P_{\rm out}^{\mathsf T}\right\|_2
   \left\|\boldsymbol P_{\rm in}\boldsymbol D_1^{-1}
   \boldsymbol D_2\boldsymbol P_{\rm in}^{\mathsf T}\right\|_2
   \right]^{-1}.
\end{align}
Then
\begin{equation}
 c_-^{12}\mathcal G_{D_2}
 \le \mathcal G_{D_1}
 \le c_+^{12}\mathcal G_{D_2}.
 \label{eq:dyn_norm_equivalence_bound}
\end{equation}

\emph{Proof.}
Coordinate selectors commute with the two diagonal scalings on their selected
subspaces.  The observed matrix in norm 1 is therefore the observed matrix in
norm 2 multiplied by the two restricted change-of-scale matrices.  Singular
value submultiplicativity gives the upper inequality; interchanging 1 and 2
gives the lower inequality. \hfill$\square$

Thus finite-time boundedness or blow-up is invariant within an equivalence
class of fixed finite-dimensional norms.  A numerical threshold-crossing time
is not invariant: it must always be reported together with
$\boldsymbol D_q$, the selectors, and the threshold.

\subsubsection{Continuous direction--wavenumber selection}

First choose the compact covering domain
\begin{equation}
 \widehat\Omega_{nk}=\mathbb S^2\times[k_{\min},k_{\max}],
 \qquad 0<k_{\min}<k_{\max}<\infty,
 \label{eq:dyn_covering_search_domain}
\end{equation}
and let
$\Omega_{nk}=\mathbb{RP}^2\times[k_{\min},k_{\max}]$ denote its antipodal
quotient.  The matrix-valued generator is defined on $\mathbb S^2$, not
directly on $\mathbb{RP}^2$; the scalar gain descends to the quotient only after
the conjugacy in Eq.~\eqref{eq:dyn_antipodal_symmetry} is established.  Define
the continuous gain envelope and the finite-time critical time by
\begin{align}
 \Gamma_D(t)
 &=\max_{([\boldsymbol n],k)\in\Omega_{nk}}
   \mathcal G_D(t;k,\boldsymbol n),
 \label{eq:dyn_gain_envelope}\\
 t_c^{\rm FT}(\boldsymbol D_q,G_c)
 &=\inf\{t\in[t_0,t_f]:\Gamma_D(t)\ge G_c\},
 \qquad G_c>1,
 \label{eq:dyn_finite_critical_time}
\end{align}
with $\inf\varnothing=+\infty$.  The registered calculation uses $G_c=e$, so
onset means one natural-log unit of maximum observable amplification in the
locked dimensionless norm; it is not a material constant.
The displayed $t_c^{\rm FT}(\boldsymbol D_q,G_c)$ is shorthand: it also depends
on $\boldsymbol P_{\rm in}$, $\boldsymbol P_{\rm out}$, $t_0$, the base path,
material parameters, and $\Omega_{nk}$.  No monotonicity of $\Gamma_D(t)$ is
assumed.

\paragraph{Proposition 4 (existence and antipodal symmetry).}
If H1--H5 hold and the admitted blocks depend continuously on
$(\boldsymbol n,k)\in\widehat\Omega_{nk}$, then the maximum in
Eq.~\eqref{eq:dyn_gain_envelope} is attained.  For real base-state tangents and
real selectors/scales,
\begin{equation}
 \boldsymbol A_{\mathrm{dyn}}(t;k,-\boldsymbol n)
 =\overline{\boldsymbol A_{\mathrm{dyn}}(t;k,\boldsymbol n)},
 \qquad
 \mathcal G_D(t;k,-\boldsymbol n)
 =\mathcal G_D(t;k,\boldsymbol n).
 \label{eq:dyn_antipodal_symmetry}
\end{equation}

\emph{Proof.}
Continuous dependence of a linear initial-value problem on its coefficients
makes $\boldsymbol\Phi$ and hence $\mathcal G_D$ continuous on
$\widehat\Omega_{nk}$.  Reversing $\boldsymbol n$ changes the
sign of every $\mathrm{i}k$ directional block and leaves the $k^2$ blocks
unchanged, which complex-conjugates the real-coefficient generator and its
propagator.  Complex-conjugate matrices have the same singular values.
Consequently, the continuous scalar gain descends to the compact quotient
$\Omega_{nk}$, where the Weierstrass theorem gives an attained maximum.
\hfill$\square$

Uniqueness is not presumed.  For a relative gain tolerance
$0<\varepsilon_G<1$, the robust selector is the near-optimal set
\begin{equation}
 \mathcal S_\varepsilon(t)
 =\left\{([\boldsymbol n],k)\in\Omega_{nk}:
 \mathcal G_D(t;k,\boldsymbol n)
 \ge(1-\varepsilon_G)\Gamma_D(t)\right\}.
 \label{eq:dyn_near_optimal_set}
\end{equation}
For fixed nonzero $\varepsilon_G$, this set generally has nonzero extent even
around a smooth strict maximum.  Numerical resolution requires Hausdorff
stability of $\mathcal S_\varepsilon$ and convergence of the maximizer candidates
under time-step, angular, radial, and search-box refinement.  A unique
$(\boldsymbol n_{\rm FT}^*,k_{\rm FT}^*)$ is reported only when the limiting
maximum is strict and isolated and the conclusion persists as
$\varepsilon_G\rightarrow0$; otherwise the selector is explicitly set-valued.
The associated linear wavelength is
$\lambda_{\rm FT}^*=2\pi/k_{\rm FT}^*$ only for an interior, refinement-stable
wavenumber.

\subsubsection{Long-wave and high-wavenumber robustness}

For $k>0$, introduce
\begin{equation}
 \widehat{\boldsymbol q}_w
 =[\widehat{\boldsymbol w},\dot{\widehat{\boldsymbol w}},
   \widehat{\boldsymbol d}]^{\mathsf T}
 =\boldsymbol T_k\widehat{\boldsymbol q},
 \qquad
 \boldsymbol T_k=\operatorname{diag}
 (\mathrm{i}k\boldsymbol I_3,\mathrm{i}k\boldsymbol I_3,
  \boldsymbol I_{63}),
 \label{eq:dyn_long_wave_transform}
\end{equation}
so that
$\boldsymbol A_w=\boldsymbol T_k\boldsymbol A_{\rm dyn}
\boldsymbol T_k^{-1}$.  If, in addition to H1--H3,
$\boldsymbol K_{\zeta\zeta}^{-1}$ remains uniform at $k=0$ and the reduced
blocks have continuous leading limits with
\begin{equation}
 \boldsymbol K_{uu}=O(k^2),\qquad
 \boldsymbol K_{ud}=O(k),\qquad
 \boldsymbol K_{du}=O(k),\qquad
 \boldsymbol K_{dd}=O(1),
 \label{eq:dyn_long_wave_block_orders}
\end{equation}
then $\boldsymbol A_w$ has a finite $k\rightarrow0^+$ limit after rigid
translation is projected out.  Because the registered input and output
selectors act only on $\boldsymbol d$, their observed propagator block is
unchanged by $\boldsymbol T_k$; the observable gain therefore has the same
finite conditional limit.  The present implementation propagates
$[\widehat{\boldsymbol u},\widehat{\boldsymbol v},\widehat{\boldsymbol d}]$
only for $k>0$ and its primary search begins at a positive $k_{\min}$; it does
not numerically certify the hypotheses in
Eq.~\eqref{eq:dyn_long_wave_block_orders}.  Consequently, a maximum that
follows $k_{\min}$ under domain enlargement is reported only as a long-wave
boundary diagnostic, not as a finite material wavelength or a proved
$k=0$ continuation.

At high wavenumber, positivity of the acoustic tensor, directional
conductivity, and the relevant gradient principal symbol is a necessary
admission test, but by itself is not a proof against non-normal transient
growth.  The following conditional statement supplies the required continuous
bound.

\paragraph{Proposition 5 (conditional uniform high-$k$ observable bound).}
Suppose that for $k\ge k_H$, all $\boldsymbol n\in\mathbb S^2$, and
$t\in I$ there is a Hermitian, differentiable symmetrizer for the fixed-scale
generator $\widetilde{\boldsymbol A}_{\rm dyn}$,
$\boldsymbol W(t;k,\boldsymbol n)>\boldsymbol0$ such that
\begin{equation}
 \dot{\boldsymbol W}
 +\boldsymbol W\widetilde{\boldsymbol A}_{\mathrm{dyn}}
 +\widetilde{\boldsymbol A}_{\mathrm{dyn}}^{\mathsf H}\boldsymbol W
 \preceq2\omega_\infty(t)\boldsymbol W,
 \qquad \omega_\infty\in L^1(I),
 \label{eq:dyn_high_k_energy_inequality}
\end{equation}
with $\omega_\infty$ independent of $k$ and $\boldsymbol n$.  Suppose also
that the selector-energy constants
\begin{align}
 C_{\rm in}
 &=\sup_{\substack{k\ge k_H\\\boldsymbol n\in\mathbb S^2}}
 \|\boldsymbol W(t_0;k,\boldsymbol n)^{1/2}
   \boldsymbol P_{\rm in}^{\mathsf T}\|_2<\infty,\\
 C_{\rm out}
 &=\sup_{\substack{t\in I,k\ge k_H\\\boldsymbol n\in\mathbb S^2}}
 \|\boldsymbol P_{\rm out}
   \boldsymbol W(t;k,\boldsymbol n)^{-1/2}\|_2<\infty
 \label{eq:dyn_high_k_selector_constants}
\end{align}
are finite.  Then
\begin{equation}
 \sup_{\substack{k\ge k_H\\\boldsymbol n\in\mathbb S^2}}
 \mathcal G_D(t;k,\boldsymbol n)
 \le
 C_{\rm out}C_{\rm in}
 \exp\!\left(\int_{t_0}^{t}\omega_\infty(\xi)\,\mathrm d\xi\right).
 \label{eq:dyn_uniform_high_k_bound}
\end{equation}

\emph{Proof.}
For $V=\widetilde{\boldsymbol q}^{\mathsf H}
\boldsymbol W\widetilde{\boldsymbol q}$,
Eq.~\eqref{eq:dyn_high_k_energy_inequality} gives
$\dot V\le2\omega_\infty V$.  The input selector enters this energy norm with
factor $C_{\rm in}$, Gronwall propagates it, and the output selector returns to
the Euclidean observation norm with factor $C_{\rm out}$, yielding
Eq.~\eqref{eq:dyn_uniform_high_k_bound}. \hfill$\square$

Equation~\eqref{eq:dyn_uniform_high_k_bound} proves uniform finite-time
boundedness of the declared observable under its stated symmetrizer and
selector hypotheses; it does not prove full-state boundedness or asymptotic
decay.  A stronger uniform equivalence
$m\boldsymbol I\preceq\boldsymbol W\preceq M\boldsymbol I$ would bound the
full scaled propagator by $\sqrt{M/m}$ times the same exponential, but standard
acoustic displacement--velocity coordinates need not admit such a $k$-uniform
equivalence.  In the present workflow, positive principal-symbol margins,
algebraic conditioning, modal residuals, and extension of the numerical scan
are reported only as consistency diagnostics along the audited paths.  They do
not construct $\boldsymbol W$, do not verify the hypothesis on a continuum,
and are not promoted to an unconditional theorem for all $k$.

\subsubsection{Dominant mechanism and nonlinear-width bridge}

Let $\widetilde{\boldsymbol q}_{\rm out}^*$ be a unit left singular vector of
the maximizing observed propagator, embedded in the 69-coordinate scaled
state.  For the mutually disjoint blocks
$b\in\{\mathrm{thermal},\mathrm{plastic},\mathrm{dislocation},
\mathrm{mechanical}\}$, define
\begin{equation}
 \eta_b^{\rm FT,out}
 =\frac{\|\boldsymbol P_b\widetilde{\boldsymbol q}_{\rm out}^*\|_2^2}
 {\sum_c\|\boldsymbol P_c\widetilde{\boldsymbol q}_{\rm out}^*\|_2^2},
 \qquad \sum_b\eta_b^{\rm FT,out}=1.
 \label{eq:dyn_mechanism_participation}
\end{equation}
Here the plastic block contains the plastic-distortion chart and signed slip;
the dislocation block contains mobile and dipole densities.  A zero mechanical
fraction follows automatically when mechanical coordinates are excluded by the
output selector and does not imply that inertia was removed from the dynamics.
The algebraic microfield may be reconstructed and reported separately.  To
make ``nearly repeated'' operational before inspecting a mechanism vector, we
define
\begin{equation}
 \delta_\sigma=\frac{\sigma_1-\sigma_2}{\sigma_1},\qquad
 \delta_\sigma<0.05
 \quad\Longleftrightarrow\quad
 \frac{\sigma_1}{\sigma_2}<1.0526 .
 \label{eq:dyn_singular_degeneracy_gate}
\end{equation}
When this gate is triggered, only the range of the leading singular subspace
is interpreted; fractions from one arbitrary basis vector are suppressed.
For a rank-one observed map, such as a scalar temperature output, $\sigma_2$
does not exist and the ratio is reported as not applicable rather than as
evidence of separation.  Onset and terminal ratios are reported for every
selector in the reproducibility table.
When $\mathcal S_\varepsilon$ contains distinct direction--wavenumber branches,
Eq.~\eqref{eq:dyn_mechanism_participation} is evaluated for every branch and
reported branchwise or as a range; one winning branch is not used to represent
the whole near-optimal set.
The $\eta_b^{\rm FT,out}$ are locked-scale coordinate fractions of an optimal singular
vector.  They are not energy fractions, dissipation fractions, or identified
causal contributions.

For nonlinear and linear windows
$I_{\rm NL}=[t_{0,\rm NL},t_{\rm NL}]$ and
$I_{\rm L}=[t_{0,\rm L},t_{\rm L}]$, the downstream computational bridge is
written
\begin{equation}
 W_{b,\mathrm{FWHM}}^{(c)}(t_{\rm NL})
 =C_{\rm FWHM}^{(c,b)}(I_{\rm NL};I_{\rm L})
  \frac{2\pi}{k_{\rm FT}^*(t_{\rm L})},
 \label{eq:dyn_fwhm_bridge}
\end{equation}
where $c$ identifies the perturbation class and $b$ the temperature or plastic
shear-rate observable.  The coefficient $C_{\rm FWHM}^{(c,b)}$ is accepted only after
nested-grid convergence with one fixed grain map.  An onset-wavelength ratio is
retained only as a memory diagnostic, and a set-valued $k_{\rm FT}^*$ induces a reported
interval rather than a single coefficient.  A same-window quantitative
relation requires both $t_{0,\rm L}=t_{0,\rm NL}$ and
$t_{\rm L}=t_{\rm NL}$.  Matching only the terminal time is labelled
same-terminal-time and remains conditioned on unequal history windows; unequal
terminal times define a cross-horizon diagnostic.  Neither is silently
relabeled as a same-window calibration.  The planned classical benchmark
uses constant top velocity, a linear initial velocity field, uniform initial
temperature, and separate unperturbed, broadband-perturbed, and
orientation-heterogeneous cases.  Equation~\eqref{eq:dyn_fwhm_bridge} is a
computational relation, not an experimental parameter identification or an
independent full-field validation claim.
